\documentclass{article}
\usepackage{inputenc}
\usepackage{fontenc}
\usepackage[ngerman]{babel}
\usepackage{amsmath,amssymb,amsthm}
\usepackage{hyperref}

\title{On Large Deviations in the ABC Conjecture and Prime Factor Constraints}
\author{Thomas Hoffbauer}
\date{}

\newtheorem{theorem}{Theorem}
\newtheorem{lemma}{Lemma}
\newtheorem{remark}{Remark}

\newcommand{\rad}{\operatorname{rad}}

\begin{document}
	\maketitle
	
	\begin{abstract}
		We study large deviations in the ABC conjecture via the quantity
		\[
		\Delta(a,b,c) := \log c - \log \rad(abc).
		\]
		Numerical experiments indicate a strong negative correlation between $\Delta$
		and the number $\omega(abc)$ of distinct prime factors.
		
		We formulate a structural hypothesis stating that large values of $\Delta$
		can only occur when the prime support is small, and discuss how this
		is consistent with known results on S-unit equations and linear forms
		in logarithms.
	\end{abstract}
	
	\section{Introduction}
	
	The ABC conjecture states that for coprime integers $a,b$ with $a+b=c$,
	and for every $\varepsilon>0$, there exists $K_\varepsilon$ such that
	\[
	c \le K_\varepsilon \cdot \rad(abc)^{1+\varepsilon}.
	\]
	
	We define
	\[
	\Delta(a,b,c) := \log c - \log \rad(abc).
	\]
	
	The conjecture implies that $\Delta$ cannot grow significantly.
	
	\section{Numerical Evidence}
	
	Computations for coprime pairs $(a,b)$ with $a,b \le N$ (up to several million)
	show:
	
	\begin{itemize}
		\item $\Delta$ is typically negative,
		\item large values of $\Delta$ are rare,
		\item large $\Delta$ occurs only when $\omega(abc)$ is small.
	\end{itemize}
	
	These observations motivate the following hypothesis.
	
	\section{Structural Hypothesis}
	
	\begin{lemma}[Structural Constraint]
		For every $\delta>0$, there exists a constant $K(\delta)$ such that
		\[
		\Delta(a,b,c) \ge \delta \implies \omega(abc) \le K(\delta).
		\]
	\end{lemma}
	
	\begin{remark}
		We do not prove this lemma; it is supported by numerical evidence
		and consistency with known diophantine constraints.
	\end{remark}
	
	\section{Connection to S-Unit Equations}
	
	Let
	\[
	S := \{p : p \mid abc\}, \quad s = |S|.
	\]
	
	Setting
	\[
	x = \frac{a}{c}, \quad y = \frac{b}{c},
	\]
	we obtain
	\[
	x + y = 1,
	\]
	where $x,y$ are S-units.
	
	By classical results (see \cite{EvertseGyory}),
	the number of such solutions is finite for fixed $S$.
	
	\section{Baker-Type Bounds}
	
	Results on linear forms in logarithms (see \cite{BakerWustholz})
	yield bounds of the form
	\[
	\log c \le C(s) \cdot P \cdot (\log P)^s,
	\]
	where $P$ is the largest prime divisor of $abc$ and $C(s)$ grows at most exponentially in $s$.
	
	\section{Heuristic Consequences}
	
	Using the prime number theorem,
	\[
	\log \rad(abc) \gtrsim s \log s.
	\]
	
	Thus, for large $s$, the radical term dominates, and $\Delta$ cannot be large.
	
	\section{Quantitative Behaviour}
	
	The above bounds suggest the existence of a function $K(\delta)$
	such that
	\[
	\Delta \ge \delta \implies \omega(abc) \le K(\delta),
	\]
	with $K(\delta)$ growing slowly as $\delta \to 0$.
	
	\section{Discussion}
	
	This perspective suggests that large deviations in the ABC conjecture
	are structurally constrained and correspond to highly restricted
	prime factor configurations.
	
	\section{Outlook}
	
	A possible route toward the ABC conjecture is to rigorously establish
	the structural constraint or derive it from known diophantine estimates.
	
	\begin{thebibliography}{9}
		
		\bibitem{BakerWustholz}
		A.~Baker and G.~W\"ustholz,
		\textit{Logarithmic Forms and Diophantine Geometry},
		Cambridge University Press, 2007.
		
		\bibitem{EvertseGyory}
		J.-H.~Evertse and K.~Gy\H{o}ry,
		\textit{Unit Equations in Diophantine Number Theory},
		Cambridge University Press, 2015.
		
	\end{thebibliography}
	
\end{document}